\par\noindent\textbf{(1)} 远场复振幅是瞳函数的二维Fourier变换，径向积分可写U(θ)∝∫\_0\textasciicircum{}a r dr∫\_0\textasciicircum{}\{2π\}exp[−ikr sinθcosφ]dφ。利用角积分等于2πJ0(kr sinθ)，并对r积分，得到U∝2πaJ1(ka sinθ)/(k sinθ)，其中Jm为第一类Bessel函数。
\par\noindent\textbf{(2)} 除以中心U(0)=πa²，得归一幅度A(x)=2J1(x)/x、x=ka sinθ，强度I/I(0)=[2J1(x)/x]²。中心极限J1(x)≈x/2使I(0)有限；将x=0直接代入形式分母为零而宣称奇异是错误的。
\par\noindent\textbf{(3)} 首暗环要求J1(x)=0的首正根x1≈3.831706，故sinθ1≈3.831706/(ka)=0.60983λ/a=1.21967λ/D，D=2a。Rayleigh判据把另一点源的Airy峰置于第一点源首暗环，得到小角分辨尺度θ\_R≈1.22λ/D。
\par\noindent\textbf{(4)} 孔径D加倍则θ\_R约减半、主峰角宽变窄，但总集光量和噪声也需考虑。有限光谱带宽令不同λ的暗环位置不一，积分后暗环可能不再严格为零；非均匀瞳面与像差亦改变理想Bessel振幅。
